\documentclass[11pt]{article}
\usepackage[margin=25mm]{geometry}
\usepackage{amsmath,amssymb,amsthm,hyperref}
\newtheorem{proposition}{Proposition}
\title{Event probabilities in the C160 moment representation}
\author{}
\date{23 October 2026}
\begin{document}
\maketitle

The first C160 block already certifies all sixteen event probabilities.
Its features are the thirty-two multilinear Bernstein monomials in the
five physical parameters. The certificate follows directly from the
two-qubit outcome amplitudes, so no semidefinite optimization is needed.

Write $z=(k,a_0,b_0,a_1,b_1)$ and define
\[
 d(t)=(1-t)^2+t^2,\qquad
 D=(1+k^2)\prod_{t\in\{a_0,b_0,a_1,b_1\}}d(t).
\]
The state is $(|00\rangle+k|11\rangle)/\sqrt{1+k^2}$.
For the observable
\[
 O(t)=\frac{1-2t}{d(t)}Z+\frac{2t(1-t)}{d(t)}X,
\]
unnormalized eigenvectors for outcomes $+1$ and $-1$ are
$f_+(t)=(1-t,t)$ and $f_-(t)=(-t,1-t)$.
They are orthogonal and each has squared norm $d(t)$.
Thus the numerator of the amplitude for outcomes $\alpha,\beta$
at settings $x,y$ is
\[
 P_{\alpha\beta xy}
 =f_{\alpha,0}(a_x)f_{\beta,0}(b_y)
   +k f_{\alpha,1}(a_x)f_{\beta,1}(b_y).
\]
Born's rule gives the following four-square identity. Here
$e_0(t)=1-t$, $e_1(t)=t$, $\bar x=1-x$ and $\bar y=1-y$:
\begin{equation}\label{eq:event}
 D\,p(\alpha,\beta|x,y)
 =\sum_{r,s=0}^1
 \bigl[P_{\alpha\beta xy}e_r(a_{\bar x})e_s(b_{\bar y})\bigr]^2.
\end{equation}
Every root is multilinear in all five variables. Consequently it has
a unique expansion in
\[
 \phi_\varepsilon(z)=\prod_{i=0}^4
 z_i^{\varepsilon_i}(1-z_i)^{1-\varepsilon_i},\qquad
 \varepsilon\in\{0,1\}^5.
\]
The coefficient of $\phi_\varepsilon$ is simply the root's value at
the cube vertex $\varepsilon$. This gives four integer vectors
$r_{\alpha\beta xy}^{rs}\in\mathbb Z^{32}$ and the explicit Gram
\[
 G_{\alpha\beta xy}
 =\sum_{r,s=0}^1r_{\alpha\beta xy}^{rs}
                  (r_{\alpha\beta xy}^{rs})^{\mathsf T}\succeq0.
\]
All sixteen Grams have rank four.

\begin{proposition}
Every point of the normalized C160 moment projection defines
nonnegative, normalized and no-signalling binary outcome probabilities.
\end{proposition}
\begin{proof}
Let $\ell$ be its moment functional, with $\ell(D)=1$, and let
$M_0(\ell)=[\ell(\phi_i\phi_j)]\succeq0$ be the first block.
The physical coordinate maps define $q$ by
$q_\nu=\ell(D\Phi_\nu)$. Applying $\ell$ to \eqref{eq:event} gives
\[
 \frac{1+\alpha q_{A_x}+\beta q_{B_y}
                  +\alpha\beta q_{E_{xy}}}{4}
 =\operatorname{tr}(G_{\alpha\beta xy}M_0(\ell))\ge0.
\]
Summing over both outcomes gives one. Summing over either outcome
gives $(1+\alpha q_{A_x})/2$ or $(1+\beta q_{B_y})/2$,
independently of the other party's setting.
\end{proof}

The sum of the four event Grams represents $D$. It need not equal
the strictly positive normalization Gram
$Q_D=\left(\begin{smallmatrix}I_{16}&I_{16}\\I_{16}&2I_{16}\end{smallmatrix}\right)$
entry by entry: Gram representations of the same polynomial are
not unique. The exact check compares their polynomial coefficients.

The file \nolinkurl{C160_EVENT_GRAMS.json} contains all sixty-four
integer root vectors. The accompanying
\nolinkurl{check_c160_event_grams.py} reconstructs them from the
amplitudes and verifies all sixteen identities against the actual
243-row C160 coefficient map, including its off-diagonal factors.
These event certificates establish probability positivity. Equality
of the C160 projection with the quantum set still requires certificates
for every quantum Bell inequality.
\end{document}
